therefore @@M0@@; now @@M1@@
@@M2@@

ii) By substitution, it suffices to prove formula (6.5.3) for @@M3@@ and for @@M4@@, @@M5@@. Since @@M6@@ is torsion-free, we can embed this ring in its tensor product with @@M7@@, and thus see that @@M8@@ is injective. Put @@M9@@, where @@M10@@ is the first nonzero term, @@M11@@, of @@M12@@. The multiplicativity of @@M13@@ gives, by (6.5.2):
@@M14@@
whence the value of @@M15@@ by identifying the terms of lowest degree.

\sgaStableTarget{sga6.e05.para.6-6}6.6. The @@M16@@-ring @@M17@@ is augmented over @@M18@@: indeed, we define a @@M19@@-homomorphism from @@M20@@ to @@M21@@ by @@M22@@. It is therefore endowed with a canonical filtration of an augmented @@M23@@-ring. We shall compare it with the one just defined.

\sgaStableTarget{sga6.e05.prop.6-6-1}\textbf{Proposition 6.6.1.} \uline{Let} @@M24@@ \uline{be an element of the form} @@M25@@. \uline{Then}:
@@M26@@
\uline{the unwritten terms} \uline{being of degree} @@M27@@.

By functoriality, it suffices to give the proof in the case
